\documentclass[11pt]{article}

\usepackage[margin=1in]{geometry}
\usepackage[T1]{fontenc}
\usepackage{lmodern}
\usepackage{microtype}
\usepackage{amsmath,amssymb,amsthm,mathtools}
\usepackage{booktabs,tabularx}
\usepackage{enumitem}
\usepackage{aliascnt}
\usepackage[dvipsnames]{xcolor}
\usepackage[colorlinks=true,linkcolor=MidnightBlue,citecolor=BrickRed,
            urlcolor=MidnightBlue]{hyperref}
\usepackage[nameinlink,capitalise,noabbrev]{cleveref}

\setlength{\parindent}{0pt}
\setlength{\parskip}{0.55em}
\setlist{itemsep=0.22em,topsep=0.35em}
\allowdisplaybreaks

\newtheorem{theorem}{Theorem}[section]
\newaliascnt{lemma}{theorem}
\newtheorem{lemma}[lemma]{Lemma}
\aliascntresetthe{lemma}
\newaliascnt{proposition}{theorem}
\newtheorem{proposition}[proposition]{Proposition}
\aliascntresetthe{proposition}
\newaliascnt{corollary}{theorem}
\newtheorem{corollary}[corollary]{Corollary}
\aliascntresetthe{corollary}
\theoremstyle{definition}
\newaliascnt{definition}{theorem}
\newtheorem{definition}[definition]{Definition}
\aliascntresetthe{definition}
\theoremstyle{remark}
\newaliascnt{remark}{theorem}
\newtheorem{remark}[remark]{Remark}
\aliascntresetthe{remark}

\newcommand{\bits}{\{0,1\}}
\newcommand{\F}{\mathbb F}
\newcommand{\C}{\mathrm C}
\newcommand{\MP}{\mathrm{MP}}
\newcommand{\Span}{\operatorname{span}}
\newcommand{\wt}{\operatorname{wt}}
\newcommand{\supp}{\operatorname{supp}}
\newcommand{\poly}{\operatorname{poly}}
\newcommand{\eps}{\varepsilon}
\newcommand{\had}{\mathbin{\odot}}
\newcommand{\status}[1]{\textsc{#1}}

% Set \anonymoustrue for an author-anonymous review copy.
\newif\ifanonymous
\anonymousfalse

\title{\textbf{Coding Witness Fibers}\\[0.35em]
\large Exact reductions, shared-evaluation tricks, and barriers for
existential circuit projection}
\ifanonymous
  \author{}
  \date{}
  \hypersetup{
    pdftitle={Coding Witness Fibers},
    pdfauthor={}
  }
\else
  \author{Samuel Schlesinger}
  \date{September 2026}
  \hypersetup{
    pdftitle={Coding Witness Fibers},
    pdfauthor={Samuel Schlesinger}
  }
\fi

\begin{document}
\maketitle

\begin{abstract}
Let $N\ge1$, $M\ge0$, let $f:\bits^N\times\bits^M\to\bits$ be a Boolean
verifier, and let
$g(x)=\bigvee_y f(x,y)$.  We ask when the $2^M$ witness branches can be
compressed into a small deterministic circuit for $g$.  A companion
mass-production theorem can batch all branches whenever the deterministic
input occupies a fixed positive fraction of the verifier input, but its
resulting bound is worse by a factor $2^M$ than direct synthesis of the
$N$-input function $g$.  The useful change of viewpoint is therefore to code
the entire witness fiber
\[
  v_x=(f(x,y))_{y\in\bits^M},
\]
not merely to batch its coordinates.

This viewpoint gives three exact reductions.  If the span of the realized
fibers has dimension $R$, an information set of $R$ actual witnesses
determines whether every fiber is zero.  If every nonempty fiber has density
at least $\delta$ and there are $K$ distinct nonempty fibers, a fixed hitting
set of about $(\log(K+1))/\delta$ actual witnesses suffices.  Without either
promise, an $f$-dependent linear map with at most $\min\{N,2^M\}$ output bits
separates every realized nonzero fiber from zero; this worst-case measurement
count is exact.  With one-sided distributional error $0<\rho\le1$, only
$\lceil\log_2(1/\rho)\rceil$ random parity measurements are needed.

The last reduction compresses measurements, not gates.  A recovery-boundary
identity does expose a shared-input trick: when parity branches are decoded
through a linear local-recovery code, repeated resource incidences cancel by
symmetric difference.  Its circuit benefit is conditional on the surviving
resource functions having a cheap joint evaluator.  We then prove an exact
obstruction: if one fixed linear sketch supports a correct update through
coordinatewise AND on all profiles, its kernel is a coordinate ideal.  Up to
a change of basis, such a sketch merely retains actual witness coordinates.
Thus a stronger construction must exploit reachable profiles, change
encodings between gates, tolerate error, or use a projection-stable
representation.  We derive concrete bounds for dense fibers, low witness
degree, and supplied decomposable or bounded-width knowledge-compilation
representations, and isolate the remaining aggregate-synthesis problem.  The
elementary coding arguments are not claimed as literature-novel; the purpose
of this note is to make the reductions, barrier, cancellation opportunity,
and open target precise.
\end{abstract}

\clearpage

\section{The question and the answer in one page}

Fix a circuit for
\[
  f:\bits^N\times\bits^M\longrightarrow\bits,
  \qquad
  g(x)=\bigvee_{y\in\bits^M}f(x,y),
  \tag{1.1}\label{eq:projection}
\]
and write $S=\C(f)$.  Throughout, $N\ge1$ and $M\ge0$; if $N=0$, the
projection is one free constant and needs no circuit.  The $x$ bits are the
deterministic input and the $y$ bits are nondeterministic choices.  The
elementary simulation uses one
restricted copy of $f$ per witness:
\[
  \C(g)\le 2^M S+2^M-1.
  \tag{1.2}\label{eq:replication-intro}
\]
The central question is whether common input and coding let us replace the
factor $2^M$ by something smaller.

The answer has two layers.  At the level of information, yes: the $2^M$
witness branches can often be replaced by a short list of actual witnesses,
and they can always be replaced by at most $N$ parity measurements when
$N\le2^M$.  At the level of unrestricted circuit gates, not automatically:
forming a parity such as
\[
  x\longmapsto\bigoplus_{y\in T}f(x,y)
  \tag{1.3}\label{eq:parity-aggregate}
\]
may itself contain essentially all of the work.

Here is the map.  Put $D=2^M$; let $K$ count the distinct nonempty witness
fibers, let $R$ be their span dimension, and let $\delta$ be their minimum
relative weight; $0<\rho\le1$ is a target distributional error.  If $K=0$,
then $g=0$ is free, so the nontrivial rows below assume $K\ge1$.

\begin{center}
\scriptsize
\begin{tabularx}{\linewidth}{@{}
  >{\raggedright\arraybackslash}p{0.17\linewidth}
  >{\raggedright\arraybackslash}p{0.27\linewidth}
  >{\raggedright\arraybackslash}p{0.20\linewidth}
  >{\raggedright\arraybackslash}X@{}}
\toprule
Method & Hypothesis and dependence & Summaries & Gate conclusion \\
\midrule
Full enumeration
  & none
  & $D$ actual witnesses
  & $D(S+1)-1$ gates \\
Span information set
  & $f$-dependent; fiber-span dimension $R$
  & $R$ actual witnesses
  & $R(S+1)-1$ gates \\
Dense-fiber hitting
  & density $\delta$; $f$-dependent
  & $O((\log(K+1))/\delta)$ actual witnesses
  & same factor times $S+1$ \\
Exact linear fingerprint
  & no promise; $f$-dependent masks
  & $O(\log(K+1))$ parities; at most $\min\{N,D\}$
  & exact zero test; aggregate cost unpaid \\
Approximate fingerprint
  & fixed $\mu$: $f,\mu$-dependent mask; or a fresh mask per evaluation
  & $\lceil\log_2(1/\rho)\rceil$ parities
  & one-sided error $\rho$; aggregate cost unpaid \\
DNNF forgetting
  & a decomposable NNF (DNNF) is supplied
  & decomposable subcircuits
  & linear-size projection from a supplied DNNF \\
\bottomrule
\end{tabularx}
\end{center}

\paragraph{The coding move.}
Regard $x\mapsto v_x$ as a codebook of at most $2^N$ realized words, and
choose an $f$-dependent linear map whose kernel avoids every realized nonzero
word.  This replaces the $D$ branch values by $O(\log(K+1))$ parity
aggregates.  It is an exact compression of the zero test, not yet a smaller
circuit: the aggregates still have to be computed from the succinct
verifier.  \Cref{thm:exact-fingerprint} proves the compression.  In
\cref{prop:recovery-boundary}, local-recovery equations cancel resource
coordinates used an even number of times, leaving a symmetric-difference
boundary whose joint evaluator is still charged.  Finally,
\cref{thm:ideal-kernel} rules out universal propagation by one fixed linear
sketch.

The most useful conceptual distinction is
\[
  \boxed{\text{few summaries}\quad\ne\quad\text{few gates to produce them}.}
  \tag{1.4}\label{eq:central-distinction}
\]
The ideal-kernel theorem in \cref{sec:and-barrier} makes this distinction
formal for fixed linear sketches and AND gates.

\clearpage
\tableofcontents
\clearpage

\section*{Scope, contribution, and research process}
\addcontentsline{toc}{section}{Scope, contribution, and research process}

\paragraph{Contribution and scope.}
The contribution claimed here is the integrated reduction-and-barrier
framework, including the recovery-boundary bridge from local recovery to
common-input parity aggregation, and the resulting specification of the
aggregate-synthesis target.  No individual elementary coding lemma is claimed
as literature-novel, and an unsuccessful search is not used as evidence of
priority.

\paragraph{Status of the claims.}
Formal statements are marked \status{proved here}, \status{imported}, or
\status{known}.  External facts and comparisons in the prose are cited where
they are used.  Conditional reductions and unresolved directions are called
conditional or open in the surrounding text.  The only imported technical
theorem is the mass-production estimate from the companion manuscript
\cite{Schlesinger2026}.

\paragraph{Research-process disclosure.}
This investigation was developed with substantive AI assistance.  The
underlying mass-production project began in a ChatGPT.com conversation that
the author exported as a \LaTeX{} document.  Later interactive sessions used
Claude Fable~5 and 5.1 through Claude Code and GPT-5.6 Sol through Codex
for experiments, broad literature searches, proof ideas, counterexample
searches, and drafting.  The author checked the arguments and citations and is
responsible for the note.

\section{Baselines and what mass production does not yet buy}
\label{sec:baselines}

We use fan-in-two $\{\wedge,\vee,\neg\}$ circuits, with free constants,
fan-out, and output designations; size counts gates.  For a Boolean map $h$
with any finite output length, $\C(h)$ is its minimum size.  Let
\[
  L(N)=\max_{h:\bits^N\to\bits}\C(h).
\]
In this convention, as $N\to\infty$, classical synthesis gives
\[
  L(N)\le (1+o(1))\frac{2^N}{N};
  \tag{2.1}\label{eq:lupanov}
\]
the sharp standard-basis statement used here is recorded by Frandsen and
Miltersen \cite{FrandsenMiltersen2005}, following Lupanov's synthesis method
\cite{Lupanov1958}.

For $t\ge1$, let $f^{\times t}$ evaluate $f$ on $t$ independent input blocks,
and set $\MP_t(f)=\C(f^{\times t})$.

\begin{proposition}[The unconditional envelope; \status{proved here}]
\label{prop:envelope}
For every $N,M$ and $f$ as in \eqref{eq:projection},
\[
  \C(g)\le
  \min\left\{
    2^M\C(f)+2^M-1,
    L(N),
    \MP_{2^M}(f)+2^M-1
  \right\}.
  \tag{2.2}\label{eq:envelope}
\]
\end{proposition}

\begin{proof}
For the first bound, hardwire a different witness into each of $2^M$ copies
of a circuit for $f$, identify all copies of the $x$ inputs, and join the
outputs with a binary OR tree.  For the second, synthesize the $N$-input
function $g$ directly.  For the third, begin with a circuit for
$f^{\times 2^M}$, identify its deterministic input blocks, hardwire its
witness blocks to all distinct values, and append the OR tree.  Restrictions,
input identification, and repeated output designations add no gates.
\end{proof}

\subsection{The imported mass-production estimate}

The companion manuscript proves the following independent-input statement.

\begin{theorem}[Exponential-range mass production; \status{imported}]
\label{thm:imported-mp}
For every fixed $0\le\gamma<1$, there is $A_\gamma$ such that every
$h:\bits^n\to\bits$ and every $1\le t\le2^{\gamma n}$ satisfy
\[
  \MP_t(h)\le A_\gamma\frac{2^n}{n}.
  \tag{2.3}\label{eq:imported-mp}
\]
The high-rate version has normalized leading coefficient at most
$1/(1-\gamma)$ \cite{Schlesinger2026}.
\end{theorem}

Apply this with $n=N+M$ and $t=2^M$.  One fixed $\gamma<1$ suffices along a
family precisely when
\[
  \frac{M}{N+M}\le\gamma,
  \tag{2.4}\label{eq:fixed-ratio}
\]
so the deterministic fraction remains bounded away from zero.  Along a
sequence with $N+M\to\infty$ and $M/(N+M)\to\theta<1$, first use any fixed
$\gamma>\theta$ and then let $\gamma\downarrow\theta$.  Since
$(1-\theta)(N+M)=N+o(N+M)$, the imported coefficient gives
\[
  \C(g)
  \le (1+o(1))\frac{2^{N+M}}{N}.
  \tag{2.5}\label{eq:mp-projection}
\]
The appended OR tree is lower order because
$2^M/(2^{N+M}/N)=N/2^N=o(1)$.  But direct synthesis gives
$(1+o(1))2^N/N$.  The mass-production route is therefore worse by the factor
$(1+o(1))2^M$.  This compares upper bounds, not the true complexity of a
particular $g$.

The conclusion is useful because it prevents a false start: the companion
theorem does batch all witnesses in the fixed-ratio regime, but that fact
alone gives no new worst-case determinization bound.  Existential projection
has reduced the exposed input length from $N+M$ to $N$, and generic synthesis
already exploits that reduction.

\subsection{What identical deterministic inputs do remove}

The common $x$ is not irrelevant.  It removes repeated evaluations of the
same resource function at the same input.

\begin{lemma}[Identical-call collapse; \status{proved here}]
\label{lem:identical-call}
Suppose a multi-output computation requests functions
$H_{i_1},\ldots,H_{i_s}:\bits^N\to\bits$ all on the same input $x$.  Replacing
the list by one representative of every distinct extensional function changes
neither minimum circuit size nor the obtainable outputs, apart from free
output designations.
\end{lemma}

\begin{proof}
From representatives to the original list, repeat each representative output
wire wherever it is requested.  In the other direction, discard duplicate
outputs.  Both operations are free in the circuit model.
\end{proof}

For existential projection, every branch supplies the same $x$.  Duplicate
requests for the same restricted function of $x$ can therefore be coalesced
by free fan-out, but up to $2^M$ distinct restricted functions may remain.
Common input compresses \emph{call multiplicity}; it does not by itself
compress \emph{function multiplicity}.  The stronger parity-specific
cancellation available inside a linear recovery code is isolated in
\cref{sec:recovery-cancellation}.

\section{The fiber code: dimension and distance from zero}
\label{sec:fiber-code}

Put $D=2^M$ and index the coordinates of $\F_2^D$ by witnesses.  For each
$x\in\bits^N$, define its witness fiber vector
\[
  v_x=(f(x,y))_{y\in\bits^M}\in\F_2^D.
  \tag{3.1}\label{eq:fiber-vector}
\]
Then $g(x)=1$ if and only if $v_x\ne0$.  Let
\[
  \mathcal V_f=\{v_x:v_x\ne0\},\qquad
  K=|\mathcal V_f|,
  \qquad
  U_f=\Span(\mathcal V_f),\qquad
  R=\dim U_f.
  \tag{3.2}\label{eq:fiber-parameters}
\]
Thus $U_f$ is a binary linear code of length $D$.  The realized nonzero
fibers need not include all of its nonzero codewords.

\subsection{Dimension gives actual witnesses}

\begin{lemma}[Information set; \status{proved here}]
\label{lem:information-set}
If $U\le\F_2^D$ has dimension $R$, there is a coordinate set
$H\subseteq[D]$ of size $R$ such that restriction $U\to\F_2^H$ is injective.
\end{lemma}

\begin{proof}
Put a basis of $U$ in the rows of an $R\times D$ matrix.  The matrix has $R$
linearly independent columns.  Their indices form $H$, and restriction to
those columns maps the chosen basis to an invertible $R\times R$ matrix.
\end{proof}

Translating coordinate indices back to witnesses gives an immediate circuit
bound.

\begin{corollary}[Span-dimension determinization; \status{proved here}]
\label{cor:span-bound}
If $R>0$, there is a fixed set $H\subseteq\bits^M$ of $R$ witnesses such that
\[
  g(x)=\bigvee_{y\in H}f(x,y),
  \qquad
  \C(g)\le R\C(f)+R-1.
  \tag{3.3}\label{eq:span-bound}
\]
If $R=0$, then $g$ is constant zero.
\end{corollary}

This is a genuine gate bound: the summaries are actual restricted copies of
the verifier.  It is nonuniform because the information set may depend on the
complete relation.  Given the explicit fibers, Gaussian elimination finds it;
no efficient procedure from a succinct circuit for $f$ is implied.

\subsection{Realized-fiber density gives a hitting set}

For a realized fiber, write $F_x=\{y:f(x,y)=1\}$.  Suppose every nonempty
$F_x$ has relative size at least $\delta$.  The number of distinct nonempty
sets is $K$, not necessarily $2^N$.

\begin{theorem}[Dense-fiber hitting; \status{proved here}]
\label{thm:dense-hitting}
If $K\ge1$ and $0<\delta<1$, there is a fixed witness set $H$ with
\[
 |H|\le
 \min\left\{
   D,
   \left\lfloor\frac{\ln K}{-\ln(1-\delta)}\right\rfloor+1
 \right\}
 \le \min\left\{D,\frac{\ln K}{\delta}+1\right\}.
 \tag{3.4}\label{eq:dense-hitting}
\]
It intersects every nonempty fiber.  For $\delta=1$, one witness suffices.
Consequently,
\[
  \C(g)\le |H|(\C(f)+1)-1.
  \tag{3.5}\label{eq:dense-circuit}
\]
\end{theorem}

\begin{proof}
Choose $h$ uniform independent witnesses.  A fixed fiber is missed with
probability at most $(1-\delta)^h$.  A union bound over the $K$ distinct
fibers makes the failure probability less than one for the displayed value of
$h$.  Remove repetitions, or use all $D$ witnesses if that is shorter.  The
second inequality uses $-\ln(1-\delta)\ge\delta$.  Finally, hardwire the
witnesses in $H$ into $|H|$ verifier copies and OR their outputs.
\end{proof}

The same proof has a greedy interpretation: among the fibers not yet hit,
some witness lies in at least a $\delta$ fraction.  This constructs $H$ from
an explicit relation, but the required coverage queries may themselves be
exponential from a succinct verifier.

If the set system $\{F_x\}$ has VC dimension $v$, put
$\bar v=\max\{1,v\}$.  The epsilon-net theorem of Haussler and Welzl
\cite{HausslerWelzl1987}, with its endpoint cases absorbed explicitly, gives
the boundary-safe alternative size
\[
  O\!\left(\frac{\bar v}{\delta}
    \left(1+\log\frac1\delta\right)\right)
  \tag{3.6}\label{eq:vc-net}
\]
under the uniform measure on the witness cube.  The measure and the common
density promise are essential.  This refinement is useful only when its
expression beats both \eqref{eq:replication-intro} and direct synthesis.

Dimension and distance are complementary.  An information set works for
every nonzero vector in the entire span $U_f$ and costs $R$ coordinates.  The
density argument needs to hit only the $K$ realized fibers and may cost much
less than $R$; it uses their actual weights, which can exceed the minimum
distance of $U_f$.

\subsection{Low witness degree}

Low degree is a canonical concrete promise under which the density route
becomes effective.  Every Boolean function of $y$ has a unique multilinear
polynomial over $\F_2$.  Suppose that, for every $x$, the polynomial
$y\mapsto f(x,y)$ has degree at most an integer $0\le d\le M$.  Its evaluation
vector lies in the binary Reed--Muller code $\operatorname{RM}(d,M)$
\cite{Muller1954,Reed1954}.

\begin{lemma}[Reed--Muller weight; \status{proved here}]
\label{lem:rm-weight}
A nonzero multilinear polynomial in $M$ variables of degree at most $d$ is
one on at least $2^{M-d}$ points.
\end{lemma}

\begin{proof}
The case $M=0$ is the nonzero constant polynomial of weight one.  For $M>0$,
induct on $M$.  Write
$p(y_1,\ldots,y_M)=a(y_1,\ldots,y_{M-1})+y_Mb(y_1,\ldots,y_{M-1})$.
If $b=0$, the weight doubles and induction applies to $a$.  If $b\ne0$, then
$\deg b\le d-1$.  Whenever $b(z)=1$, exactly one of $p(z,0)$ and $p(z,1)$ is
one.  Induction gives at least $2^{(M-1)-(d-1)}=2^{M-d}$ such $z$.
\end{proof}

Thus $\delta\ge2^{-d}$ in \cref{thm:dense-hitting}, giving an
$f$-dependent witness set of size $O(2^d(\log K+1))$.  There is also an
explicit universal information set:
\[
  H_{M,d}=\{y\in\bits^M:\wt(y)\le d\},
  \qquad
  |H_{M,d}|=B(M,d):=\sum_{i=0}^d\binom Mi.
  \tag{3.7}\label{eq:low-weight-set}
\]

\begin{lemma}[Low-weight information set; \status{proved here}]
\label{lem:low-weight-info}
If a degree-at-most-$d$ multilinear polynomial vanishes on $H_{M,d}$, it is
zero.  Moreover, no smaller coordinate set works for the entire degree class.
\end{lemma}

\begin{proof}
Write $p(y)=\sum_{|A|\le d}c_A\prod_{i\in A}y_i$.  At the indicator $1_S$,
where $|S|\le d$,
\[
  p(1_S)=\sum_{A\subseteq S}c_A.
\]
Induction on $|S|$ recovers every $c_S$ from these values.  Hence restriction
to $H_{M,d}$ is injective.  The polynomial space has dimension $B(M,d)$, so
rank--nullity forbids an injective restriction to fewer coordinates.
\end{proof}

If $K=0$, then $g$ is constant zero and $\C(g)=0$.  Assume $K\ge1$ and define
the density term
\[
  Q_d(K)=
  \begin{cases}
    1,&d=0,\\[0.2em]
    \left\lfloor\dfrac{\ln K}{-\ln(1-2^{-d})}\right\rfloor+1,
      &1\le d\le M.
  \end{cases}
\]
Combining the three actual-witness routes gives the exact-form bound
\[
  \C(g)\le L_f(\C(f)+1)-1,
  \qquad
  L_f=\min\{D,R,B(M,d),Q_d(K)\},
  \tag{3.8}\label{eq:low-degree-combined}
\]
whenever the degree promise holds.  For fixed $d$,
$B(M,d)=\Theta(M^d)$, while the relation-dependent density branch is
$O(2^d\log(K+1))$.

\section{Zero-preserving parity fingerprints}
\label{sec:fingerprints}

Span dimension and density can both be unfavorable.  Measurement compression
is nevertheless always possible because only finitely many fibers are
realized.

\begin{theorem}[Exact linear fingerprint; \status{proved here}]
\label{thm:exact-fingerprint}
If $K=0$, the zero map suffices.  If $K\ge1$, put
\[
  r=\min\left\{D,\max\{1,\lceil\log_2K\rceil\}\right\}.
  \tag{4.1}\label{eq:fingerprint-width}
\]
There is an $\F_2$-linear map $A:\F_2^D\to\F_2^r$ such that
\[
  Av_x=0\quad\Longleftrightarrow\quad v_x=0
  \qquad\text{for every }x\in\bits^N.
  \tag{4.2}\label{eq:zero-preserving}
\]
For $N\ge1$, one always has $r\le\min\{N,D\}$, and this worst-case bound is
exact over all relations $f$.
\end{theorem}

\begin{proof}
For the upper bound, maintain the nonzero fibers not yet separated from zero
by the chosen rows.  A uniform linear functional vanishes on each fixed
nonzero vector with probability $1/2$, so some row leaves at most half of the
current vectors unseparated.  After $r-1$ rows, at most two remain.  A final
functional can be one on a single remaining vector.  It can also be one on
two, because two distinct nonzero vectors over $\F_2$ are linearly
independent.  This separates the last vectors.  The inequality $K\le2^N$
gives $r\le N$, and
$K\le2^D-1$ gives $r\le D$.

For tightness, first suppose $D\le N$.  A relation can realize every vector
of $\F_2^D$ among its fibers, forcing any successful $A$ to be injective and
to have at least $D$ output bits.  If $D>N$, realize zero and every nonzero
vector of an $N$-dimensional subspace, using exactly $2^N$ deterministic
inputs.  The restriction of $A$ to that subspace must be injective, so it
needs at least $N$ output bits.
\end{proof}

The construction is deterministic given the explicit fiber list: at each
step choose a row attaining the averaging bound.  It is not claimed to be
efficient from a succinct circuit for $f$.

The zero-collision argument is a specialization of standard linear-hashing
ideas.  Universal and linear hashing study broader collision and load goals
\cite{CarterWegman1979,AlonEtAl1999}; succinct perfect linear hashing can even
separate every pair in a prescribed sparse set \cite{ChenJinWilliams2019}.
Here only the pairs $(0,v)$ for $v\in\mathcal V_f$ must be separated, and no
efficient hash construction from the succinct verifier is asserted.

For the rest of this section assume $K\ge1$, so $r\ge1$; the case $K=0$ has
$g=0$ and needs no circuit.

Write the rows of $A$ as $a_1,\ldots,a_r$ and define
\[
  h_i(x)=\langle a_i,v_x\rangle
  =\bigoplus_{y:a_i(y)=1}f(x,y).
  \tag{4.3}\label{eq:fingerprint-aggregates}
\]
Then
\[
  g(x)=\bigvee_{i=1}^r h_i(x).
  \tag{4.4}\label{eq:fingerprint-zero-test}
\]
Equivalently, after choosing a basis of $\F_{2^r}$, the columns of $A$ are
weights $\alpha_y\in\F_{2^r}$ and
\[
  g(x)=1
  \quad\Longleftrightarrow\quad
  \sum_y\alpha_y f(x,y)\ne0.
  \tag{4.5}\label{eq:field-fingerprint}
\]
One field symbol still consists of $r$ Boolean wires; the notation does not
make the sum cheap.

\subsection{Width depends only on the target error}

The exact map must work on all $x$ at once.  For distributional approximation
of $g$, the width loses its dependence on $N$.

\begin{theorem}[Approximate parity fingerprint; \status{proved here}]
\label{thm:approx-fingerprint}
Let $\mu$ be any fixed distribution on $\bits^N$ and let $t\ge0$ be an
integer.  Choose $t$ independent uniform rows and define
\[
  \widetilde g_A(x)=1[Av_x\ne0].
\]
For every $A$, $\widetilde g_A$ has no false positives.  Moreover,
\[
  \mathbb E_A\Pr_{x\sim\mu}
  [g(x)=1,\widetilde g_A(x)=0]
  =2^{-t}\Pr_{x\sim\mu}[g(x)=1]
  \le2^{-t}.
  \tag{4.6}\label{eq:approx-fingerprint}
\]
Consequently some fixed, $f$- and $\mu$-dependent $A$ has one-sided error at
most $2^{-t}$.  A freshly random $A$ has the same error bound pointwise for
every fixed input.
\end{theorem}

\begin{proof}
Zero always maps to zero.  For $v_x\ne0$, one uniform row has uniform inner
product with $v_x$, so all $t$ rows vanish with probability $2^{-t}$.
Interchange the finite expectations over $A$ and $x$.
\end{proof}

Thus, for $0<\rho\le1$, taking
$t=\lceil\log_2(1/\rho)\rceil$ gives false-negative mass at most $\rho$.
This is the standard randomized parity sketch of OR, applied only to the
realized witness fibers; linear-sketch and parity-query treatments place it
in a broader framework \cite{KannanEtAl2018,Gavinsky2025}.  On the full cube
$\F_2^D$, an exact deterministic linear sketch for OR must be injective and
needs $D$ rows.  \Cref{thm:exact-fingerprint} escapes that lower bound because
only at most $2^N$ fiber vectors are realized.

Small-bias sample spaces can shorten the random seed while making each
nonzero parity nearly unbiased \cite{NaorNaor1993}.  A constant-distance
linear code gives another formulation: encode $v$, sample codeword
coordinates, and use the distance to bound the probability that every sample
is zero.  These ideas reduce randomness or measurement count, not necessarily
the Boolean gates used to produce a measurement from an implicit $v_x$.

The distinction from Valiant--Vazirani isolation is also important
\cite{ValiantVazirani1986}.  Isolation intersects a witness set with random
constraints and seeks a unique surviving \emph{witness}.  Here a row reports
the parity of a set and may cancel many accepting witnesses.  The exact
quantifier order is
\[
  \forall f\;\exists A\;\forall x,
\]
not a randomized per-instance isolation claim.

\subsection{Recovery-boundary cancellation on a common input}
\label{sec:recovery-cancellation}

Linear local recovery offers a sharper common-input trick than duplicate-call
collapse.  Parity aggregates combine recovery equations over $\F_2$, so a
resource incidence used twice disappears.

\begin{lemma}[Symmetric-difference recovery; \status{proved here}]
\label{lem:symmetric-difference}
Let $E:\F_2^D\to\F_2^Q$ be an injective linear encoding, and let
$E^*:\F_2^Q\to\F_2^D$ be its adjoint under the standard inner products.  For
each message coordinate $i$, choose a recovery vector $b_i$ satisfying
$E^*b_i=e_i$, and put $R_i=\supp(b_i)$.  Then every $T\subseteq[D]$ satisfies
\[
  \bigoplus_{i\in T}v_i
  =\bigoplus_{z\in\triangle_{i\in T}R_i}(Ev)_z
  \qquad(v\in\F_2^D),
  \tag{4.7}\label{eq:symmetric-difference}
\]
where $\triangle$ denotes symmetric difference.
\end{lemma}

\begin{proof}
Injectivity of $E$ makes $E^*$ surjective, so the stated recovery vectors
exist.  Set $b_T=\sum_{i\in T}b_i$.  Then
\[
  \left\langle b_T,Ev\right\rangle
  =\left\langle E^*b_T,v\right\rangle
  =\left\langle\sum_{i\in T}e_i,v\right\rangle.
\]
Over $\F_2$, the support of $b_T$ is precisely the symmetric difference of
the supports of the $b_i$.
\end{proof}

The affine-line recovery used in the companion construction gives a vivid
special case \cite[Section~4]{Schlesinger2026}.  In characteristic two, a
degree-at-most-$(q-2)$ polynomial on an affine line $L$ of size $q$ satisfies
$P(u)=\sum_{z\in L\setminus\{u\}}P(z)$.  If every $u\in T\subseteq L$ is
recovered through that same line, then
\[
  \triangle_{u\in T}(L\setminus\{u\})
  =
  \begin{cases}
    T,&|T|\text{ even},\\
    L\setminus T,&|T|\text{ odd}.
  \end{cases}
  \tag{4.8}\label{eq:line-cancellation}
\]
In particular, the parity of $q-1$ target values collapses to the one
remaining code symbol.  The identity applies bitwise after expanding field
symbols in an $\F_2$ basis.

The next statement records exactly when this cancellation becomes a circuit
bound.  For $Z\subseteq[Q]$, let
\[
  G_Z(x)=\bigl((Ev_x)_z\bigr)_{z\in Z},
  \qquad B_E(Z)=\C(G_Z).
\]

\begin{proposition}[Recovery-boundary reduction; \status{proved here}]
\label{prop:recovery-boundary}
Let the rows $a_1,\ldots,a_r$ of $A$ separate the realized fibers from zero.
For each $j$, choose $b_j\in\F_2^Q$ with $E^*b_j=a_j$, and let
$Z=\bigcup_j\supp(b_j)$.  Then, in the standard circuit basis,
\[
  \C(g)\le B_E(Z)
  +4\sum_{j=1}^r\max\{0,\wt(b_j)-1\}+r-1.
  \tag{4.9}\label{eq:recovery-boundary-bound}
\]
\end{proposition}

\begin{proof}
Compute the distinct resource functions in $G_Z$ once.  The adjoint identity
gives
\[
  \langle a_j,v_x\rangle
  =\langle b_j,Ev_x\rangle
  =\bigoplus_{z\in\supp(b_j)}G_z(x).
\]
A zero-support parity is the free constant zero.  A parity of $s\ge1$ bits
uses $s-1$ XORs, each implementable by at most four
$\{\wedge,\vee,\neg\}$ gates.  An OR tree on the $r$ resulting fingerprint
bits uses $r-1$ further gates.
\end{proof}

The proposition suggests jointly choosing the fingerprint rows and their
recovery representations to minimize a symmetric-difference boundary.  It is
specific to the shared input: for independent inputs, two incidences with the
same resource coordinate generally carry different values and cannot cancel.
It is also only a conditional gain.  Defining an encoding whose first
coordinates are the desired aggregates makes every $b_j$ sparse
tautologically; unless those encoded coordinates inherit a cheap evaluator,
$B_E(Z)$ contains the original difficulty.

\subsection{Why this is not yet a circuit upper bound}

Suppose the $r$ functions in \eqref{eq:fingerprint-aggregates} have a joint
circuit of size $B$.  Appending an OR tree proves
\[
  \C(g)\le B+r-1.
  \tag{4.10}\label{eq:aggregate-reduction}
\]
This conditional reduction is exact.  The missing theorem is a useful bound
on $B$ in terms of $S,N,M$, and $r$.

Materializing every $f(x,y)$ first costs at most $D\C(f)$ gates, and direct
parity trees add $O(rD)$ more.  A short seed describing the masks does not
remove these gates.  Likewise, applying mass production to materialize all
fiber coordinates returns to the bound in \cref{sec:baselines}.  The new
question is whether selected aggregates can be computed \emph{without}
expanding the full fiber.

\section{Why a fixed sketch cannot simply pass through AND}
\label{sec:and-barrier}

The witness profiles form the algebra
\[
  \mathcal R=\F_2^D
\]
with coordinatewise multiplication
$(u\had v)(y)=u(y)v(y)$.  For a wire $w$ of the verifier, let
\[
  V_w(x)=(w(x,y))_{y\in\bits^M}.
\]
At an AND gate $w=u\wedge v$,
$V_w(x)=V_u(x)\had V_v(x)$.  We may replace OR by
$u\vee v=u\oplus v\oplus(u\wedge v)$ at constant cost.  In the resulting
$\{\wedge,\oplus,\neg\}$ description, XOR and NOT are linear or affine on
profiles, so AND is the nonlinear obstruction to propagating one fixed linear
fingerprint.  Recall that a subspace $I\subseteq\mathcal R$ is an ideal when
$u\had v\in I$ for every $u\in I$ and every $v\in\mathcal R$.

\begin{theorem}[Ideal-kernel classification; \status{proved here}]
\label{thm:ideal-kernel}
Let $A:\mathcal R\to\F_2^r$ be linear.  The following are equivalent.
\begin{enumerate}[label=(\roman*)]
\item There is an arbitrary function $\Phi$ on $A(\mathcal R)^2$ such that
\[
  A(u\had v)=\Phi(Au,Av)
  \qquad\text{for all }u,v\in\mathcal R.
  \tag{5.1}\label{eq:and-update}
\]
\item $\ker A$ is an ideal of the coordinatewise algebra $\mathcal R$.
\end{enumerate}
Every ideal has the form
\[
  I_T=\Span\{e_y:y\in T\}
  =\{u:\supp(u)\subseteq T\}
  \tag{5.2}\label{eq:coordinate-ideal}
\]
for a unique set of witness coordinates $T$.
\end{theorem}

\begin{proof}
If \eqref{eq:and-update} holds, $k\in\ker A$, and $v\in\mathcal R$, then
\[
  A(k\had v)=\Phi(0,Av)=A(0\had v)=0,
\]
so the kernel is an ideal.  Conversely, if $I=\ker A$ is an ideal and
$Au=Au'$, $Av=Av'$, then $u'=u+k$ and $v'=v+\ell$ for $k,\ell\in I$.  The
difference between $u'\had v'$ and $u\had v$ is
\[
  k\had v+u\had\ell+k\had\ell\in I.
\]
Hence $A(u\had v)$ depends only on $Au,Av$, and $\Phi$ is well defined.

For the classification, if $u\in I$ and $u(y)=1$, then
$e_y=e_y\had u\in I$.  Thus $I$ is spanned by exactly the unit vectors that
it contains, proving \eqref{eq:coordinate-ideal}.
\end{proof}

\begin{corollary}[Universal propagation equals witness retention;
\status{proved here}]
\label{cor:projection-barrier}
Assume $K\ge1$.  Suppose $A$ satisfies \cref{thm:ideal-kernel} and also separates every
$v\in\mathcal V_f$ from zero.  Put $H=[D]\setminus T$, where
$\ker A=I_T$.  Then $H$ intersects every nonempty witness fiber,
$\operatorname{rank}A=|H|$, and, up to an injective linear change of output
coordinates, $A$ is exactly the projection retaining the bits indexed by
$H$.  In particular,
\[
  \C(g)\le |H|(\C(f)+1)-1.
  \tag{5.3}\label{eq:barrier-hitting}
\]
\end{corollary}

\begin{proof}
The quotient $\mathcal R/I_T$ is the coordinate space on $H$.  Since $A$ has
kernel $I_T$, its induced map from that quotient to $A(\mathcal R)$ is an
isomorphism, and its rank is $|H|$.  If a realized nonzero fiber vanished on
$H$, it would lie in $I_T$ and be killed by $A$, a contradiction.  Thus $H$
hits every nonempty fiber, and the circuit bound follows by restriction and
replication.
\end{proof}

For the smallest example, take $D=2$ and $A(a,b)=a+b$.  The pairs
\[
  (u,v)=((1,0),(1,0)),\qquad
  (u',v')=((1,0),(0,1))
\]
have the same two input sketches, both equal to one, but the product sketches
are one and zero.  No update rule can distinguish them.

\paragraph{What the theorem does not rule out.}
Its quantifiers range over every $u,v\in\mathcal R$, it uses the same sketch
at both inputs and the output, and it demands exactness.  A verifier-specific
construction may need updates only on profile pairs reached by a common $x$;
it may use a different code at every wire; or it may tolerate error.  Those
are real escape routes.  The theorem rules out only the tempting universal
shortcut of pushing one short mixed-parity sketch through arbitrary ANDs.

\section{Structural escapes and their limits}
\label{sec:structural-escapes}

\subsection{Changing profile spaces: route and dominance check}

A gatewise code must pay for its transitions, not just its state names.
Suppose each wire $w$ has an encoding
$q_w:\bits^N\to\bits^{r_w}$ of its realized witness profile, each gate has a
transition circuit of size $\tau_w$, and a final circuit of size
$\tau_{\mathrm{out}}$ tests whether the output profile is zero.  Composing
the transition circuits in the verifier DAG gives
\[
  \C(g)\le \tau_{\mathrm{out}}+\sum_w\tau_w
  +\text{leaf-encoding cost}.
  \tag{6.1}\label{eq:profile-system}
\]
This is bookkeeping, but it states the open target correctly: short profile
names, cheap transitions, and an exact zero test are all required.

Output-span dimension by itself will not improve the nonuniform gate bound.
If the realized output profiles span a positive-dimensional space of dimension
$k_{\mathrm{out}}$, \cref{lem:information-set} already gives
$k_{\mathrm{out}}$ actual witnesses and a circuit of size
$k_{\mathrm{out}}(S+1)-1$; if the dimension is zero, then $g=0$.  The dense
simulation below is retained because it isolates transition cost and may
matter when coefficient tables are sparse or algorithmically available.

One sufficient certificate uses changing linear spaces.  For subspaces
$E,F\le\mathcal R$, let
$E\had F=\Span\{u\had v:u\in E,v\in F\}$.  Assign a subspace $E_w$ to every
wire, containing every realized profile, such that NOT includes the constant
one and its child space, while an AND output contains the product of its child
spaces.

\begin{proposition}[Profile-space simulation; \status{proved here}]
\label{prop:profile-space}
For a supplied AND/NOT verifier with $T$ gates, put $k_w=\dim E_w$.  Over the
mixed basis $\{\wedge,\vee,\oplus,\neg\}$, its projection has size at most
\[
  \sum_{w=\neg u} k_w(k_u+1)
  +\sum_{w=u\wedge v}(k_w+1)k_uk_v
  +\max\{0,k_{\mathrm{out}}-1\}.
  \tag{6.2}\label{eq:profile-space-bound}
\]
The standard basis changes this by at most a factor four.  In particular,
if all dimensions are at most $k$, the bound is $O(Tk^3+k)$.  A standard-basis
size-$S$ circuit can be converted to an AND/NOT circuit at constant cost, so an
$O(Sk^3+k)$ conclusion follows when an admissible dimension-$k$ family is
supplied for that converted circuit.
\end{proposition}

\begin{proof}
Choose a basis of each $E_w$ and maintain the coefficient vector of
$V_w(x)$.  A deterministic input has profile $x_i\mathbf1$ and a witness
input has a fixed profile, so their coefficients use only input wires and
constants.  NOT is an affine linear map and costs at most $k_w(k_u+1)$.
At an AND gate, expand every product of one child basis vector with another in
the output basis.  Compute the $k_uk_v$ products of input coefficients once;
the $k_w$ output coordinates are fixed parities of these products.  This gives
the second sum.  Finally, the output profile is nonzero precisely when its
coefficient vector is nonzero, and the last term is an OR tree.  XOR has a
four-gate implementation in the standard basis.
\end{proof}

Profile propagation therefore becomes interesting only when the bilinear
coefficient tables used by the gate transitions are already available or
unusually sparse, when uniform algorithmic access matters, or when one uses
nonlinear codes tailored to the realized state pairs.

Multiplication-friendly codes suggest language for the last possibility:
structured linear encodings can support reconstruction of products
\cite{CascudoEtAl2012}.  Existing constructions normally expand a
lower-dimensional algebra into enough coordinates for multiplication.  They
do not compress the full algebra $\F_2^{2^M}$.  Any application here must
control reachable subspaces or change representations between gates; otherwise
\cref{thm:ideal-kernel} applies.

\subsection{DNNF and bounded width}

There is a clean established regime in which existential projection really is
cheap.  A decomposable negation normal form (DNNF) is a negation normal form in
which the two children of every AND gate mention disjoint variable sets.

\begin{proposition}[DNNF forgetting; \status{known}]
\label{prop:dnnf}
If a DNNF of edge size $T$ for $f(x,y)$ is supplied, all $y$ variables can be
existentially forgotten in linear time without asymptotically increasing
size.  The resulting DNNF computes $g$ \cite{Darwiche1999,Darwiche2001}.
\end{proposition}

The operation replaces every positive or negative literal of a forgotten
variable by true and simplifies.  Existential quantification distributes over
OR.  At a decomposable AND, the forgotten variables of the children are
disjoint, so their witnesses can be chosen independently and quantification
also distributes.  Without decomposability, the rule is false: the formula
$y\wedge\neg y$ would incorrectly become true.

The phrase ``a DNNF is supplied'' is essential: the cited compilation theorem
does not give a polynomial-size DNNF for an arbitrary circuit.  A concrete
positive case is a CNF $F$ whose interaction graph has treewidth $w$; it
compiles to DNNF in size and time $2^{O(w)}|F|$ under the source's edge-size
convention \cite{Darwiche1999,Darwiche2001}.  If one insists that projection
preserve a complete structured deterministic DNNF, width can grow
exponentially;
Capelli and Mengel give width-parameterized upper and lower bounds
\cite{CapelliMengel2019}.  These results make the missing invariant tangible:
decomposition or separator width tells the circuit how summaries combine.

Deterministic, smooth DNNF also supports bottom-up model counting and weighted
model counting \cite{ChaviraDarwiche2008}.  That is stronger than the existence
test needed here.
Counting and parity aggregates provide useful algorithmic routes only when the
representation makes their sum-product evaluation valid; generic exact
counting hardness is context, not an unrestricted circuit lower bound
\cite{Valiant1979}.

\section{Approximation and the projection metric}
\label{sec:approximation}

An approximant to the verifier need not approximate its existential
projection.  This matters because approximation in the companion manuscript
is measured on the full $(x,y)$ cube.

\begin{proposition}[Projection can amplify error; \status{proved here}]
\label{prop:error-amplification}
Let $f'$ differ from $f$ on an $\eps$ fraction of uniformly random pairs
$(x,y)$, and let $g'(x)=\bigvee_yf'(x,y)$.  Then
\[
  \Pr_x[g(x)\ne g'(x)]\le\min\{1,2^M\eps\},
  \tag{7.1}\label{eq:error-amplification}
\]
and the factor $2^M$ is tight.
\end{proposition}

\begin{proof}
Every $x$ on which the projections differ contributes at least one erroneous
pair $(x,y)$.  Divide the number of erroneous pairs by $2^N$.  For tightness,
take $f(x,y)=1[y=0^M]$ and $f'=0$.  Their distance is $2^{-M}$, while their
projections are the constants one and zero.
\end{proof}

There are two correct one-sided alternatives.

\begin{itemize}
\item The parity fingerprint in \cref{thm:approx-fingerprint} controls error
directly over $x$, independently of fiber density, but still needs an
aggregate evaluator.
\item If nonempty fibers have density $\delta$, sample $t$ actual witnesses
and OR their verifier values.  A fixed yes-input is missed with probability at
most $(1-\delta)^t$, there are no false positives, and the circuit uses at
most $tS+t-1$ gates.  Averaging fixes one sample with the same error bound
under any chosen distribution on $x$.
\end{itemize}

Uniform $(x,y)$ approximation, direct approximation of the projection, and a
per-fiber density guarantee are therefore different notions.  Only the last
two control the decision problem in \eqref{eq:projection} without the
$2^M$ amplification.

\section{Circuit-complexity consequences}
\label{sec:consequences}

It is useful to state the parameter comparisons without complexity-class
shorthand.  If $K=0$, then $g=0$.  Otherwise, let $B_{\mathrm{agg}}$ denote
the cost of a supplied joint circuit for the $r$ parity aggregates.  Up to
lower-order and fixed-basis factors, the available envelope is
\[
\C(g)\le\min\left\{
  2^M(S+1),
  \frac{2^N}{N},
  R(S+1),
  \frac{\log(K+1)}{\delta}(S+1),
  B(M,d)(S+1),
  B_{\mathrm{agg}}+r
\right\},
\tag{8.1}\label{eq:final-envelope}
\]
where each structural term is present only under its stated hypothesis.  The
imported mass-production term is valid in the fixed-ratio regime but is
omitted from \eqref{eq:final-envelope} because direct synthesis generically
dominates it.

The replication term beats direct synthesis when roughly
\[
  M+\log_2(S+1)<N-\log_2N.
  \tag{8.2}\label{eq:replication-threshold}
\]
The dense-fiber term is a genuine improvement only if its number of retained
witnesses is below both $2^M$ and $2^N/(N(S+1))$.  Similar comparisons apply
to $R$ and $B(M,d)$.  A reduction in witnesses is not automatically a
reduction below the best circuit already available for $g$.

\subsection{Why a universal polynomial determinizer would be major}

Suppose there were a constant $c$ such that every relation with
$M\le\poly(N)$ satisfied
\[
  \C(\exists y\,f)\le(N+M+\C(f)+1)^c.
  \tag{8.3}\label{eq:hypothetical-determinizer}
\]
Here $\mathrm{NP/poly}$ is the class recognized by polynomial-size
nonuniform verifier circuits with polynomially many witness bits.  Applying
the hypothesis to those verifiers would give
$\mathrm{NP/poly}\subseteq\mathrm{P/poly}$; the reverse inclusion is trivial,
so equality follows, and in particular
$\mathrm{NP}\subseteq\mathrm{P/poly}$.  The Karp--Lipton theorem would then
collapse the polynomial hierarchy to its second level
\cite{KarpLipton1980}.  A stronger uniform compiler that takes a
polynomial-size verifier description, runs in polynomial time, and outputs a
polynomial-size circuit for its projection would imply
$\mathrm P=\mathrm{NP}$.  These are conditional consequences, not results of
this note.

The aggregate barrier has parallel context.  For a uniform verifier family,
$\bigoplus_yf(x,y)$ is a parity-counting predicate, while the integer sum is a
$\#\mathrm P$ computation.  For polynomial-size nonuniform verifier families,
a polynomial circuit bound for every parity aggregate would give
$\mathord\oplus\mathrm{P/poly}=\mathrm{P/poly}$.  Uniform counting hardness
does not itself prove a Boolean circuit lower bound, so the rigorous
unconditional obstruction remains \cref{thm:ideal-kernel}.

Unambiguity does not make the aggregate free.  If every fiber has size at
most one, OR equals parity exactly, but evaluating that parity may still be
hard.  Valiant--Vazirani's isolation theorem reinforces the distinction
between obtaining a unique-witness promise and deterministically evaluating
the resulting language \cite{ValiantVazirani1986}.

For comparison, easy-witness results compress the \emph{description of a
witness} under circuit assumptions; they do not supply one circuit computing
the projection for every $x$ \cite{MurrayWilliams2020}.  Work on
nondeterministic circuit lower bounds also addresses a different direction.
For example, Morizumi gives restricted-basis lower bounds and an explicit
Tseitin-style translation from circuits to nondeterministic 3-CNFs
\cite{Morizumi2015}.  None of those results turns the measurement reduction
in \cref{thm:exact-fingerprint} into a generic gate bound.

There is, however, an elementary lower-bound transfer.  If $g=\exists_yf$ and
$\C(g)\ge L$, then \eqref{eq:replication-intro} implies
\[
  \C(f)\ge\frac{L+1}{2^M}-1.
  \tag{8.4}\label{eq:lower-transfer}
\]
The factor $2^M$ makes this weak when the witness budget is large.

\section{The resulting research program}
\label{sec:open}

The mass-production viewpoint does inspire a technical target for this
approach, but not a finished generic determinization theorem.  It says to
regard all branches with the same deterministic input as one coded object and
to share the work that produces summaries of that object.  The preceding
results narrow the target to four forms.

\begin{enumerate}[label=\textbf{\arabic*.},leftmargin=2.2em]
\item \textbf{Boundary-efficient coded recovery.}
Choose a zero-preserving fingerprint together with recovery representations
whose symmetric-difference boundary $Z$ is small, as in
\cref{prop:recovery-boundary}, and prove that the surviving resource map
$G_Z$ has a cheap joint circuit.  The objective is the full right side of
\eqref{eq:recovery-boundary-bound}, not just the number of fingerprint bits.

\item \textbf{Changing multiplicative fingerprints.}
Use different encodings at different wires and require transitions only on
profiles reached by the verifier.  The fixed-sketch ideal theorem no longer
applies, but every transition must have an explicit small circuit.  Restricted
product spaces, sparse bilinear transition tables, and
multiplication-friendly codes are plausible tools.

\item \textbf{Succinct actual-witness hitting.}
The dense and low-degree arguments give small hitting sets, but the
relation-dependent choices inspect the fibers.  Find checkable promises under
which a small $H$ can be constructed from a succinct verifier without
enumerating its truth table.

\item \textbf{Restricted-model lower bounds.}
Exhibit verifier families for which every zero-preserving aggregate system has
large transition cost, even though its output width is small.  The
ideal-kernel theorem resolves only the fixed linear/full-algebra model; lower
bounds for changing or nonlinear sketches are not resolved here.
\end{enumerate}

The central compression hierarchy is therefore
\[
\boxed{
\begin{array}{c}
\text{fiber-code dimension or density from zero}\Rightarrow
\text{few actual witnesses}\Rightarrow\text{a circuit bound},\\[0.35em]
\text{few realized fibers}\Rightarrow
\text{few parity measurements}\not\Rightarrow\text{a circuit bound},\\[0.35em]
\text{small recovery boundary plus cheap shared resources}
\Rightarrow\text{a circuit bound}.
\end{array}}
\tag{9.1}\label{eq:hierarchy}
\]
The second implication becomes useful precisely when a structural
representation supplies the missing aggregate computation.  That is the
classical nondeterministic-sharing problem left by this note.

\section*{Reproducibility note}

The accompanying research corpus contains an exhaustive finite checker for
small instances of
\cref{thm:exact-fingerprint,lem:rm-weight,lem:low-weight-info,lem:information-set,lem:symmetric-difference,thm:ideal-kernel},
including affine-line cancellation and the two-bit AND collision.  These
checks are evidence against transcription and
edge-case errors; they are not substitutes for the proofs above.  The source
registry separately records the primary-source boundary for every imported or
known result.

\begin{thebibliography}{99}

\bibitem{AlonEtAl1999}
Noga Alon, Martin Dietzfelbinger, Peter Bro Miltersen, Erez Petrank, and
G\'abor Tardos.
\newblock Linear hash functions.
\newblock \emph{Journal of the ACM}, 46(5):667--683, 1999.
\newblock \href{https://doi.org/10.1145/324133.324179}{doi:10.1145/324133.324179}.

\bibitem{CapelliMengel2019}
Florent Capelli and Stefan Mengel.
\newblock Tractable QBF by knowledge compilation.
\newblock In \emph{36th Symposium on Theoretical Aspects of Computer Science},
LIPIcs 126, pages 18:1--18:16, 2019.
\newblock \href{https://doi.org/10.4230/LIPIcs.STACS.2019.18}{doi:10.4230/LIPIcs.STACS.2019.18}.

\bibitem{CarterWegman1979}
J. Lawrence Carter and Mark N. Wegman.
\newblock Universal classes of hash functions.
\newblock \emph{Journal of Computer and System Sciences}, 18(2):143--154,
1979.
\newblock \href{https://doi.org/10.1016/0022-0000(79)90044-8}{doi:10.1016/0022-0000(79)90044-8}.

\bibitem{CascudoEtAl2012}
Ignacio Cascudo, Ronald Cramer, Chaoping Xing, and An Yang.
\newblock Asymptotic bound for multiplication complexity in the extensions of
small finite fields.
\newblock \emph{IEEE Transactions on Information Theory}, 58(7):4930--4935,
2012.
\newblock \href{https://doi.org/10.1109/TIT.2011.2180696}{doi:10.1109/TIT.2011.2180696}.

\bibitem{ChaviraDarwiche2008}
Mark Chavira and Adnan Darwiche.
\newblock On probabilistic inference by weighted model counting.
\newblock \emph{Artificial Intelligence}, 172(6--7):772--799, 2008.
\newblock \href{https://doi.org/10.1016/j.artint.2007.11.002}{doi:10.1016/j.artint.2007.11.002}.

\bibitem{ChenJinWilliams2019}
Lijie Chen, Ce Jin, and R. Ryan Williams.
\newblock Hardness magnification for all sparse NP languages.
\newblock In \emph{60th IEEE Symposium on Foundations of Computer Science},
pages 1240--1255, 2019.
\newblock \href{https://doi.org/10.1109/FOCS.2019.00077}{doi:10.1109/FOCS.2019.00077}.

\bibitem{Darwiche1999}
Adnan Darwiche.
\newblock Compiling knowledge into decomposable negation normal form.
\newblock In \emph{Proceedings of IJCAI 1999}, pages 284--289, 1999.
\newblock \url{https://www.ijcai.org/Proceedings/99-1/Papers/042.pdf}.

\bibitem{Darwiche2001}
Adnan Darwiche.
\newblock Decomposable negation normal form.
\newblock \emph{Journal of the ACM}, 48(4):608--647, 2001.
\newblock \href{https://doi.org/10.1145/502090.502091}{doi:10.1145/502090.502091}.

\bibitem{FrandsenMiltersen2005}
Gudmund Skovbjerg Frandsen and Peter Bro Miltersen.
\newblock Reviewing bounds on the circuit size of the hardest functions.
\newblock \emph{Information Processing Letters}, 95(2):354--357, 2005.
\newblock \href{https://doi.org/10.1016/j.ipl.2005.03.009}{doi:10.1016/j.ipl.2005.03.009}.

\bibitem{Gavinsky2025}
Dmytro Gavinsky.
\newblock Unambiguous parity-query complexity.
\newblock \emph{Random Structures \& Algorithms}, 66(3):e70010, 2025.
\newblock \href{https://doi.org/10.1002/rsa.70010}{doi:10.1002/rsa.70010}.

\bibitem{HausslerWelzl1987}
David Haussler and Emo Welzl.
\newblock Epsilon-nets and simplex range queries.
\newblock \emph{Discrete \& Computational Geometry}, 2:127--151, 1987.
\newblock \href{https://doi.org/10.1007/BF02187876}{doi:10.1007/BF02187876}.

\bibitem{KannanEtAl2018}
Sampath Kannan, Elchanan Mossel, Swagato Sanyal, and Grigory Yaroslavtsev.
\newblock Linear sketching over $\F_2$.
\newblock In \emph{33rd Computational Complexity Conference}, LIPIcs 102,
pages 8:1--8:37, 2018.
\newblock \href{https://doi.org/10.4230/LIPIcs.CCC.2018.8}{doi:10.4230/LIPIcs.CCC.2018.8}.

\bibitem{KarpLipton1980}
Richard M. Karp and Richard J. Lipton.
\newblock Some connections between nonuniform and uniform complexity classes.
\newblock In \emph{Proceedings of the 12th Annual ACM Symposium on Theory of
Computing}, pages 302--309, 1980.
\newblock \href{https://doi.org/10.1145/800141.804678}{doi:10.1145/800141.804678}.

\bibitem{Lupanov1958}
Oleg B. Lupanov.
\newblock On a method of circuit synthesis.
\newblock \emph{Izvestiya VUZ, Radiofizika}, 1(1):120--140, 1958.
\newblock In Russian. \url{https://radiophysics.unn.ru/issues/1958/1/120}.

\bibitem{Morizumi2015}
Hiroki Morizumi.
\newblock Lower bounds for the size of nondeterministic circuits.
\newblock In \emph{Computing and Combinatorics}, LNCS 9198, pages 289--296,
2015.
\newblock \href{https://doi.org/10.1007/978-3-319-21398-9_23}{doi:10.1007/978-3-319-21398-9\_23}.

\bibitem{Muller1954}
David E. Muller.
\newblock Application of Boolean algebra to switching circuit design and to
error detection.
\newblock \emph{Transactions of the I.R.E. Professional Group on Electronic
Computers}, EC-3(3):6--12, 1954.
\newblock \href{https://doi.org/10.1109/IREPGELC.1954.6499441}{doi:10.1109/IREPGELC.1954.6499441}.

\bibitem{MurrayWilliams2020}
Cody D. Murray and R. Ryan Williams.
\newblock Circuit lower bounds for nondeterministic quasi-polytime from a new
easy witness lemma.
\newblock \emph{SIAM Journal on Computing}, 49:STOC18-300--STOC18-322, 2020.
\newblock \href{https://doi.org/10.1137/18M1195887}{doi:10.1137/18M1195887}.

\bibitem{NaorNaor1993}
Joseph Naor and Moni Naor.
\newblock Small-bias probability spaces: Efficient constructions and
applications.
\newblock \emph{SIAM Journal on Computing}, 22(4):838--856, 1993.
\newblock \href{https://doi.org/10.1137/0222053}{doi:10.1137/0222053}.

\bibitem{Reed1954}
Irving S. Reed.
\newblock A class of multiple-error-correcting codes and the decoding scheme.
\newblock \emph{Transactions of the IRE Professional Group on Information
Theory}, 4(4):38--49, 1954.
\newblock \href{https://doi.org/10.1109/TIT.1954.1057465}{doi:10.1109/TIT.1954.1057465}.

\bibitem{Schlesinger2026}
Samuel Schlesinger.
\newblock Exponential-range mass production of Boolean functions: Local
recovery, disjoint scheduling, and bounded congestion.
\newblock Unpublished companion manuscript, version September 3, 2026.
Distributed alongside this note.

\bibitem{Valiant1979}
Leslie G. Valiant.
\newblock The complexity of enumeration and reliability problems.
\newblock \emph{SIAM Journal on Computing}, 8(3):410--421, 1979.
\newblock \href{https://doi.org/10.1137/0208032}{doi:10.1137/0208032}.

\bibitem{ValiantVazirani1986}
Leslie G. Valiant and Vijay V. Vazirani.
\newblock NP is as easy as detecting unique solutions.
\newblock \emph{Theoretical Computer Science}, 47:85--93, 1986.
\newblock \href{https://doi.org/10.1016/0304-3975(86)90135-0}{doi:10.1016/0304-3975(86)90135-0}.

\end{thebibliography}

\end{document}
